\documentclass[12pt]{article}
\usepackage[utf8]{inputenc}
\usepackage[T1]{fontenc}
\usepackage{hyperref}
\title{Loup A Language Resources for K-12 Curriculum Development}
\author{Compiled by Claude Sonnet 5.5 (claude-sonnet-5-5) for Andr\'as Kornai}
\date{}
\begin{document}
\maketitle

\section{Introduction}
Loup A is an extinct Eastern Algonquian (Algic) language of colonial New England. It is attested in mid-eighteenth-century French missionary manuscripts, chiefly the notebook \emph{Mots loups} by Jean-Claude Mathevet, edited by Gordon Day in 1975. The data were recorded from refugees at the St.\ Francis mission in Quebec. Scholars disagree about who spoke it: one view, supported by Gustafson's grammar (2000), links it to the Nipmuck of central Massachusetts, while Day linked it to the Pocumtuck. The language has no native speakers. Revitalization work would depend on the sparse manuscript record and on relation to neighbouring Algonquian languages. I did not search separately for such revitalization work.

\section{Primary Educational Resources}
No primers, textbooks or school programmes were found in search results.

\section{Linguistic and Reference Materials}
\begin{itemize}
\item Gustafson, \emph{A Grammar of the Nipmuck Language} (University of Manitoba thesis, record page). \url{https://mspace.lib.umanitoba.ca/handle/1993/2306}
\item The same thesis, alternate record. \url{https://mspace.lib.umanitoba.ca/items/0fb2392b-6422-49c0-8498-7f45d1f1a243}
\item The same thesis, PDF hosted on nipmuck.org. \url{https://www.nipmuck.org/uploads/5/0/7/7/50775337/gustafson_thesis_nipmuck_grammar.pdf}
\item Document on nipmuck.org titled ``loup (a) sappir''. I did not verify its content. \url{https://www.nipmuck.org/uploads/5/0/7/7/50775337/loup_(a)_sappir.pdf}
\item Wikipedia, ``Loup languages'': an overview of Loup A and Loup B. \url{https://en.wikipedia.org/wiki/Loup_languages}
\end{itemize}

\section{Religious and Community Resources}
The nipmuck.org site hosts the Nipmuck-related documents above, so a Nipmuck community presence exists there. I did not verify any programme, contact or revitalization activity.

\section{Digital and Audio Resources}
Languagehat, ``Five lost languages rediscovered in Massachusetts'': a blog post. \url{https://languagehat.com/five-lost-languages-rediscovered-in-massachusetts/}
No audio resources were found.

\section{Curriculum Development Context}
The record is a single word-list-style manuscript source and a grammar reconstructed from it. Its attribution to the Nipmuck is contested. Any curriculum should be developed with the relevant Indigenous communities and should state the uncertainty. It could use the Gustafson grammar and the Day edition of the source as its base.

\end{document}